% !PRINTED 241
\weberpage{241}

214. \textit{Uebernommene Formen.} Die Formen $(f, f)^2$, $(f, (f, f)^2)$,
$(f, f)^4$, welche in diesem Systeme auftreten, haben wir bereits auch
im System der Form vierten Grades vorgefunden, und es ist die Frage,
ob sich unter den Covarianten der Form fünften Grades nicht auch
eine Covariante

\weq{1}{j = (ab)^2 (ac)^2 (bc)^2 a_x b_x c_x}

befindet, welche der Invariante

\weq{2}{j = (ab)^2 (ac)^2 (bc)^2}

der Form vierten Grades entspricht. Erinnern wir uns, dass die
Invariante $j$ unsymbolisch die Form hatte:

\weq{3}{j = 6 \begin{vmatrix} \overline{a_0} & \overline{a_1} & \overline{a_2} \\ \overline{a_1} & \overline{a_2} & \overline{a_3} \\ \overline{a_2} & \overline{a_3} & \overline{a_4} \end{vmatrix},}

so können wir nunmehr für die Form fünften Grades die analoge
Determinante bilden:

\weq{4}{j = 6 \begin{vmatrix} \overline{a_0} & \overline{a_1} & \overline{a_2} & \overline{a_3} \\ \overline{a_1} & \overline{a_2} & \overline{a_3} & \overline{a_4} \\ \overline{a_2} & \overline{a_3} & \overline{a_4} & \overline{a_5} \\ -x_2^3 & x_1 x_2^2 & -x_1^2 x_2 & x_1^3 \end{vmatrix}.}

Setzt man hierin statt der wirklichen Coefficienten die symbolischen,
so erhält man in derselben Weise, wie wir früher die Determinante
(3) in (2) umgeformt hatten, für die Determinante (4) das symbolische
Product (1).

Diese Covariante $j$ befindet sich aber nicht unmittelbar im
System, und es ist die nächste Frage, welche Stellung sie zu ihm
einnimmt. Zu dem Zwecke werden wir sie so umformen, dass sie
den Factor $(ac)^4$ erhält, was nach Nr. 98 stets möglich ist. Es ist:

\[ a_x (bc) = b_x (ac) - c_x (ab), \]

also

\[ j = (ab)^2 (ac)^3 (bc) b_x^2 c_x - (ab)^3 (ac)^2 (bc) b_x c_x^2. \]

Vertauscht man im zweiten Terme $b$ mit $c$, so kommt:

\[ j = 2 (ab)^2 (ac)^3 (bc) b_x^2 c_x. \]

Vertauscht man hierin $a$ mit $c$ und nimmt die halbe Summe der so
entstehenden Ausdrücke, so erhalten wir endlich

\[ j = -(ac)^4 (ab) (cb) b_x^3. \]

Das symbolische Product rechts wird aber durch die Ueberschiebung
erhalten:

\[ ((ac)^4 a_x c_x, b_x^5)^2 = (i, f)^2, \]
%CONT>>


% ===== pdf idx 256, printed 242.0 =====
% !PRINTED 242
\weberpage{242}

%<<CONT
d. h. $j$ ist bis auf den Factor $-1$ nichts anderes, als die bereits im
Systeme aufgenommene Ueberschiebung $(f, i)^2$.

Wir erkennen also die bemerkenswerthe Thatsache, dass sämmtliche
Formen

\weq{1}{f, \ \Delta, \ t, \ i, \ j}

des vollständigen Systemes vierter Ordnung ihre entsprechenden
Formen im System von $f = a_x^5$ besitzen, wobei die gegenseitige Beziehung
 durch die je zwei solchen Formen gemeinschaftlichen charakteristischen
 Klammerfactoren hergestellt ist. Wir nennen diese Formen

\weq{2}{f, \ (f, f), \ (f, (f, f)^2), \ (f, f)^4, \ -(f, i)^2}

\gesperrt{übernommene Formen}. Sie gehen aus dem Systeme (1) einfach
dadurch hervor, dass man die Formen desselben jedem Symbole entsprechend
 mit je einem Factor erster Art multiplicirt. Diese merkwürdige
 Thatsache wiederholt sich indessen später nicht mehr; das
System der Form fünften Grades kann nicht mehr in das System der
Form sechsten Grades in dieser Weise übernommen werden.

215. \textit{Zweite Definition der Form $j$.} Es möge gleich an dieser
Stelle eine die Form $j$ geradezu definirende Eigenschaft derselben
erwähnt werden. Wenn man nämlich die Ueberschiebungen von $f$ mit
$j$ bildet, so zeigt sich, dass die dritte Ueberschiebung $(f, j)^3$ identisch
verschwindet. Denn man hat:

\[ (f, j)^3 = -(ab)^3 (bi)^2 a_x^2. \]

Setzt man hierin

\[ a_x (bi) = b_x (ai) - i_x (ab), \]

so kommt:

\[ -(ab)^3 (bi) (ai) a_x b_x + (ab)^4 (bi) a_x i_x = (f, j)^3. \]

Beide Terme links verschwinden; der erste, weil er bei Vertauschung
von $a$ mit $b$ nur sein Zeichen ändert; der zweite, weil er nichts anderes
 darstellt, als die erste Ueberschiebung der quadratischen Form
$i = (ab)^4 a_x b_x$ über sich selbst. Es ist also in der That:

\[ (f, j)^3 = 0. \]

Durch diese Eigenschaft ist aber $j$ vollständig charakterisirt, d. h.
ist $\psi$ eine beliebige cubische Form, so beschaffen, dass

\[ (f, \psi)^3 = 0, \]

dann ist stets:

\[ \psi = j, \]

bis auf einen constanten Factor.

Ist nämlich

\[ (f, \psi)^3 = a_x^2 (a\psi)^3 = a_x^2 (a_1 \psi_2 - \psi_1 a_2)^3 = 0 \]
%CONT>>


% ===== pdf idx 257, printed 243.0 =====
% !PRINTED 243
\weberpage{243}

%<<CONT
so müssen die Coefficienten von $x$ einzeln verschwinden, d. h.
die Gleichungen bestehen:

\[
\begin{aligned}
a_1^2 (a_1 \psi_2 - \psi_1 a_2)^3 &= \overline{a_0}\, \overline{\psi_3} - 3 \overline{a_1}\, \overline{\psi_2} + 3 \overline{a_2}\, \overline{\psi_1} - \overline{a_3}\, \overline{\psi_0} = 0 \\
a_1 a_2 (a_1 \psi_2 - \psi_1 a_2)^3 &= \overline{a_1}\, \overline{\psi_3} - 3 \overline{a_2}\, \overline{\psi_2} + 3 \overline{a_3}\, \overline{\psi_1} - \overline{a_4}\, \overline{\psi_0} = 0 \\
a_2^2 (a_1 \psi_2 - \psi_1 a_2)^3 &= \overline{a_2}\, \psi_3 - 3 \overline{a_3}\, \overline{\psi_2} + 3 \overline{a_4}\, \overline{\psi_1} - \overline{a_5}\, \overline{\psi_0} = 0. % [sic] third row prints $\psi_3$ without the bar carried by the other rows
\end{aligned}
\]

Bestimmt man hieraus die Verhältnisse der $\psi_i$ und trägt deren
Werthe in

\[ \psi = \overline{\psi_0} x_1^3 + 3 \overline{\psi_1} x_1^2 x_2 + 3 \overline{\psi_2} x_1 x_2^2 + \overline{\psi_3} x_2^3 \]

ein, so erhält man gerade die Covariante $j$.

\textit{Anmerkung.} Alle Formen, welche den Factor $(aj)^3$ besitzen, verschwinden
 identisch, denn sie können durch Ueberschiebung mit

\[ (f, j)^3 = (aj)^3 a_x^2 \]

entstanden gedacht werden. Der Klammerfactor $(aj)^3$ ist demnach
\gesperrt{Reducent}.

216. \textit{Ersetzung der Formen des Systemes durch eines ihrer Glieder.}
Während nun die übernommenen Formen bereits einen einfachen symbolischen
 Ausdruck besitzen, der für die weiteren Rechnungen tauglich
 ist, sind die meisten der übrigen Formen des Systemes durch
Aggregat symbolischer Producte dargestellt, das die symbolische Rechnung
 sehr erschweren würde. Es ist daher angezeigt, hier von dem
Satze Gebrauch zu machen, dass sich jede Ueberschiebung durch eines
ihrer Glieder und niedrigere Ueberschiebungen darstellen lässt. Indem
wir alsdann von den niedrigeren Ueberschiebungen, als bekannten
Bildungen, absehen, können wir direct das betreffende Glied, das ja
alle charakteristischen Merkmale der Ueberschiebung besitzt, welcher
es angehört, an Stelle der Ueberschiebung selbst in das Formensystem
einfügen. Insbesondere wollen wir nun im Folgenden die quadratischen
und linearen Covarianten, sowie die Invarianten durch eines ihrer
Glieder ersetzen, soweit die betreffenden Ueberschiebungen nicht ohnehin
 schon durch ein einfaches symbolisches Product dargestellt sind.

217. \textit{Die drei quadratischen Covarianten.} Als erste quadratische
Covariante erhielten wir:

\[ i = i_x^2 = (ab)^4 a_x b_x. \]

Zu einer zweiten giebt die Hesse'sche Form $\tau$ der cubischen Covariante
$j$ Veranlassung:

\[ \tau = \tau_x^2 = (j j_1)^2 j_x j_{1x}. \]

Diese Form $\tau$ lässt sich auch noch in anderer Weise darstellen. Es ist:

\weq{1}{(f, (j, i))^3 = (a_x^5, (ji) j_x^2 i_x)^3 = a_x^2 (aj)^2 (ai) (ji).}
%CONT>>


% ===== pdf idx 258, printed 244.0 =====
% !PRINTED 244
\weberpage{244}

%<<CONT
Setzt man nach dem Identitätssatze:

\weq{2}{a_x (ji) = -i_x (aj) + j_x (ai)}

und führt diesen Werth von $a_x (ji)$ in (1) ein, so kommt, da der
erste Term wegen des Reducenten $(aj)^3$ null ist:

\[ (f, (ji))^3 = a_x j_x (aj)^2 (ai)^2. \]

Der Ausdruck rechts kann als zweite Ueberschiebung von

\[ a_x^3 (ai)^2 \quad \text{über} \quad j_x^3 \]

angesehen werden, oder auch, weil $a_x^3 (ai)^2 = -j_x^3$, als zweite Ueberschiebung
 von $-j_x^3$ über $j_x^3$. Es ist also:

\[ (f, (j, i))^3 = -(j j_1)^2 j_x j_{1x} = -\tau. \]

Die dritte quadratische Covariante endlich, die wir mit in unsere
nächsten Betrachtungen hereinziehen, ist die Functionaldeterminante
von $i$ und $\tau$, nämlich

\[ \vartheta = (i, \tau) = (i \tau) i_x \tau_x. \]

Es fragt sich nun, in welchem Zusammenhange $\tau$ und $\vartheta$ mit den
beiden andern im Systeme auftretenden quadratischen Covarianten
$(\varphi, i^2)^4$ und $(\varphi, i^3)^5$ stehen.

Was die Covariante $\tau$ betrifft, so erkennt man sie leicht als ein
Glied der Ueberschiebung $(\varphi, i^2)^4$. Denn da $j = -(ai)^2 a_x^3$, so ist

\[ \tau = (j j')^2 j_x j'_x = ((ai)^2 a_x^3, (bi')^2 b_x^3)^2 = (ab)^2 (ai)^2 (bi')^2 a_x b_x. \]

Der Term rechts kann aber durch viermalige Faltung des Productes

\[ \varphi \cdot i^2 = (ab)^2 a_x^3 b_x^3 \cdot i_x^2 \cdot i_x'^2 \]

erhalten werden, wodurch die Behauptung erwiesen ist.

Ganz ebenso ist $\vartheta$ ein Glied der Covariante $(\varphi, i^3)^5$. Denn

\[ (i, \tau) = (i_x''^2, (ab)^2 (ai)^2 (bi')^2 a_x b_x) = -(ab)^2 (ai)^2 (bi')^2 (ai'') b_x i_x''. \]

Der Ausdruck rechts geht durch fünfmalige Faltung aus dem Product

\[ (ab)^2 a_x^3 b_x^3 \cdot i_x^2 \cdot i_x'^2 \cdot i_x''^2 \]

hervor, ist also in der That ein Glied von $(\varphi, i^3)^5$.

218. \textit{Die fundamentalen Invarianten.} Wir hatten im Systeme zunächst
 drei gerade fundamentale Invarianten erhalten; es waren dies
die Formen

\[ (i, i)^2, \quad (\varphi, i^3)^6, \quad (f^2, i^5)^{10}. \]

Die erste derselben behalten wir im Folgenden bei und bezeichnen
sie mit

\[ A_{ii} = A = (i, i)^2. \]

An Stelle der zweiten setzen wir die zweite Ueberschiebung von $i$
über $\tau$, d. i. die Invariante

\[ B = (i, \tau)^2. \]


% ===== pdf idx 259, printed 245.0 =====
% !PRINTED 245
\weberpage{245}
Sie ist gleichfalls vom achten Grade in den Co\"efficienten und man erkennt sie als ein Glied der Ueberschiebung $(\varphi,\ i^3)^6$. Denn es ist
\[
(i,\ \tau)^2 = (i_x^{\prime\prime 2},\ (ab)^2(ai)^2(bi')^2 a_x b_x)^2 = (ab)^2(ai)^2(bi')^2(ai'')(bi''),
\]
ein Term, der auch durch sechsmalige Faltung aus $(ab)^2 a_x^3 b_x^3 \cdot i_x^2 i_x^{\prime 2} i_x^{\prime\prime 2}$ erhalten wird. Die dritte Invariante des Systemes ersetzen wir durch die Discriminante von $j$, oder (was das n\"amliche ist) durch die Discriminante von $\tau$; wir erhalten
\[
(\tau,\ \tau)^2 = C.
\]
Ihren Zusammenhang mit $(f^2, i^5)^{10}$ werden wir sp\"ater kennen lernen. (Vgl. Nr. 226). Da $\tau$ vom sechsten Grade ist in den Co\"efficienten von $f$, so ist $C$ vom zw\"olften Grade wie $(f^2, i^5)^{10}$.

219. \textit{Die vier linearen Covarianten.} Wir hatten im System als einfachste lineare Covariante erhalten:
\weq{1}{(f,\ i^2)^4 = (a_x^5,\ i_x^2 \cdot i_x^{\prime 2})^4 = (ai)^2(ai')^2 a_x\,.}
Wir bezeichnen sie im Folgenden mit $\alpha$, und man erkennt aus dem Factor $(ai)^2$, dass $\alpha$ betrachtet werden kann als zweite Ueberschiebung von $j = -(ai)^2 a_x^3$ \"uber $i$, d.\,h. man hat:
\[
\alpha = -(ji)^2 j_x\,.
\]
Als zweite lineare Covariante trat im Systeme die Form $(f^3, i)^5$ auf. Wir ersetzen sie durch
\weq{2}{\beta = (i,\ \alpha)\,.}
Beide Formen stimmen \"uberein, denn man hat einestheils:
\[
(i,\ \alpha) = (i_x^2,\ (i'a)^2(i''a)^2 a_x) = (ia)(i'a)^2(i''a)^2 i_x
\]
und anderntheils
\[
-(f,\ i^3)^5 = (i^3,\ f)^5 = (i_x^2 i_x^{\prime 2} i_x^{\prime\prime 2},\ a_x^5)^5 = (ia)(i'a)^2(i''a)^2 i_x\,.
\]
Die dritte lineare Covariante $\gamma$ wird von der Ueberschiebung
\[
(t,\ i^4)^8 = ((ab)^2(ca)\,a_x^2 b_x^3 c_x^4,\ i_{1x}^2 i_{2x}^2 i_{3x}^2 i_{4x}^2)^8
\]
geliefert. Wir benutzen von ihr n\"amlich das Glied
\[
\gamma = (ab)^2(ac)(ai_1)^2(bi_2)^2 b_x (ci_3)^2(ci_4)^2\,.
\]
F\"uhren wir rechts das Symbol $\alpha_x = (ci_3)^2(ci_4)^2 c_x$ ein, so geht $\gamma$ \"uber in
\[
\gamma = (ab)^2(a\alpha)(ai_1)^2(bi_2)^2 b_x\,.
\]
Das symbolische Product rechts ist aber nichts anderes als die erste Ueberschiebung von
\[
\tau = (jj')^2 j_x j'_x = (ab)^2(ai)^2(bi')^2 a_x b_x
\]
\"uber $\alpha_x$;
%CONT>>


% ===== pdf idx 260, printed 246.0 =====
% !PRINTED 246
\weberpage{246}
%<<CONT
daher k\"onnen wir die lineare Covariante $\gamma$ darstellen durch
\weq{3}{\gamma = (\tau,\ \alpha)\,.}
Was endlich die vierte lineare Covariante $(t, i^5)^9$ betrifft, welche im Systeme auftritt, so ersetzen wir sie zun\"achst durch ihr Glied:
\[
(ab)^2(ac)(ia)^2(i_1b)^2(i_2c)^2(i_3c)^2(i_4b)\,i_{4x} = G\,.
\]
Verm\"oge der Identit\"at: $(ac)i_{4x} = (i_4c)a_x - (i_4a)c_x$ geht dasselbe \"uber in:
\[
\begin{aligned}
G &= (ab)^2(ia)^2(i_1b)^2(i_2c)^2(i_3c)^2(i_4c)(i_4b)\,a_x \\
  &\quad - (ab)^2(ia)^2(i_1b)^2(i_4a)(i_4b)\cdot(i_2c)^2(i_3c)^2\,c_x = G_1 - G_2\,.
\end{aligned}
\]
Der Term $G_2$ besitzt den Factor $\alpha_x = (i_2c)^2(i_3c)^2 c_x$, w\"ahrend sein anderer Factor nichts anderes ist als die Invariante
\[
B = (\tau,\ i)^2 = ((ab)^2(ai)^2(bi_1)^2 a_x b_x,\ i_{4x}^2)^2\,.
\]
Wir k\"onnen daher von diesem Terme $G_2 = B \cdot \alpha$, als einem Producte bekannter Formen, absehen, und nunmehr $G_1$ als die einzuf\"uhrende lineare Invariante $\delta$ definiren:
\[
\delta = (ab)^2(ia)^2(i_1b)^2(i_2c)^2(i_3c)^2(i_4c)(i_4b)\,a_x\,.
\]
Ber\"ucksichtigen wir, dass $(i_2c)^2(i_3c^2)(i_4c)i_{4x} = \beta_x$, so k\"onnen wir diese Covariante $\delta$ (vom Vorzeichen abgesehen) auch darstellen durch
\[
\begin{aligned}
\delta &= (ab)^2(ia^2)(i_1b)^2(b\beta)\,a_x \\
       &= ((ab)^2(ia)^2(i_1b)^2 b_x a_x,\ \beta_x) = ((jj_1)^2 j_x j_{1x},\ \beta_x),
\end{aligned}
\]
oder endlich durch
\weq{4}{\delta = (\tau,\ \beta)\,.{}^{*)}}

220. \textit{Die schiefe Invariante.} Wir greifen aus der Ueberschiebung $(f\cdot t,\ i^7)^{14}$, welche wir als schiefe Invariante erkannt haben, jenes Glied heraus, welches durch viermalige Faltung von $a_x^5$ mit $i_1^2$, also durch Bildung von $\alpha$, sodann durch weitere achtmalige Faltung von $b_x^2 c_x^3 d_x^4$ mit $i_2^4$, also durch Bildung von $\gamma$, und endlich durch zweimalige Faltung von $\alpha \cdot \gamma$ mit $i_3$ entsteht. Dieses Glied
\[
R = (i_x^2,\ \alpha_x \gamma_x)^2 = (i_x^2,\ (\tau\alpha)\tau_x \cdot \alpha_x)^2
\]
oder
\weq{5}{R = (i\tau)(\tau\alpha)(i\alpha)}
f\"uhren wir als schiefe Invariante an Stelle von $(f\cdot t,\ i^7)^{14}$ in das System ein.

Wir k\"onnen das symbolische Product (5) noch in verschiedener Weise durch Ueberschiebung erzeugen; so ist
\weq{6}{(i\tau)(\tau\alpha)(i\alpha) = -((\tau\alpha)\tau_x,\ (i\alpha)i_x) = -(\gamma,\ \beta) = (\beta\gamma) = -(\tau\alpha)(\tau\beta)}

${}^{*)}$ \gesperrt{Clebsch} f\"uhrt in seinen \glqq Bin\"aren Formen\grqq{} als vierte lineare Covariante die Form $(\vartheta,\ \alpha)$ ein.
%CONT>>


% ===== pdf idx 261, printed 247.0 =====
% !PRINTED 247
\weberpage{247}
%<<CONT
oder
\weq{7}{(i\tau)(\tau\alpha)(i\alpha) = ((i\alpha)(i\tau)\tau_x,\ \alpha_x) = -(\delta,\ \alpha) = (\alpha\delta)}
oder
\weq{8}{(i\tau)(\tau\alpha)(i\alpha) = ((i\tau)i_x\tau_x,\ \alpha_x^2)^2 = (\vartheta,\ \alpha^2)^2\,.}

221. \textit{Tabelle der wichtigsten Formen.} Ehe wir nun zu Anwendungen \"ubergehen, wollen wir noch einmal jene Formen zusammenstellen, die wir in den folgenden Untersuchungen vorzugsweise benutzen werden. Es sind dies

1) die \"ubernommenen Formen
\[
\begin{aligned}
f &= a_x^5 \\
\varphi &= (ab)^2 a_x^3 b_x^3 \qquad t = (ab)^2(ca)\,a_x^2 b_x^3 c_x^4 \\
i &= (ab)^4 a_x b_x \qquad\ \ -j = (ai)^2 a_x^3,
\end{aligned}
\]

2) die Covarianten:
\[
\begin{aligned}
\tau &= (jj_1)^2 j_x j_{1x} = -(f,\ (j,\ i))^3 \\
\vartheta &= (i\tau)i_x\tau_x \\
\alpha &= (ai)^2(ai_1)^2 a_x = -(ji)^2 j_x \\
\beta &= (i\alpha)i_x \\
\gamma &= (\tau\alpha)\tau_x \\
\delta &= (\tau\beta)\tau_x\,,
\end{aligned}
\]

3) die Invarianten:
\[
A = (i,\ i)^2,\quad B = (i,\ \tau)^2,\quad C = (\tau,\ \tau)^2,
\]
\[
R = (\beta\gamma) = (\alpha\delta) = (\vartheta\alpha)^2 = (i\tau)(\tau\alpha)(i\alpha) = -(\tau\alpha)(\tau\beta)\,.
\]
Hiezu treten noch

4) drei Hilfsinvarianten
\[
\begin{aligned}
M &= (i\alpha)^2 = (\beta\alpha) \\
N &= (\tau\alpha)^2 = (\gamma\alpha) \\
P &= (\tau\beta)^2 = (\delta\beta),
\end{aligned}
\]
welche bezw. vom 12., 16. und 20. Grade in den Co\"efficienten sind, w\"ahrend $A$, $B$, $C$ bezw. die Grade 4, 8, 12 besitzen. Wir werden die Darstellung der Formen $M$, $N$, $P$ durch die drei Fundamentalinvarianten in den folgenden Paragraphen kennen lernen.

\begin{center}
\textbf{\S~23.}\\
\textbf{Relationen zwischen den In- und Covarianten der Form f\"unften Grades.}
\end{center}

222. \textit{Die Quadrate der linearen Covarianten $\beta$, $\gamma$, $\delta$.} Die im vorigen Paragraphen eingef\"uhrten linearen Covarianten $\beta$, $\gamma$, $\delta$ haben wir in Functionaldeterminantenform dargestellt; ihre Quadrate m\"ussen
%CONT>>


% ===== pdf idx 262, printed 248.0 =====
% !PRINTED 248
\weberpage{248}
%<<CONT
sich also durch einfachere Formen ausdr\"ucken lassen. In der That, erhebt man zun\"achst die Identit\"at
\[
(ii_1)\alpha_x = (\alpha i_1)i_x - (\alpha i)i_{1x}
\]
in das Quadrat, so erh\"alt man:
\[
(ii_1)^2 \cdot \alpha^2 = (\alpha i_1)^2 \cdot i + (\alpha i)^2 \cdot i_1 - 2(\alpha i_1)(\alpha i)i_{1x}i_x
\]
oder:
\weq{1}{A \cdot \alpha^2 = 2M \cdot i - 2\beta^2\,.}
Ebenso findet man unter Benutzung der Identit\"aten
\[
\begin{aligned}
(\tau\tau_1)\alpha &= (\alpha\tau_1)\tau_x - (\alpha\tau)\tau_{1x} \\
(\tau\tau_1)\beta &= (\beta\tau_1)\tau_x - (\beta\tau)\tau_{1x}
\end{aligned}
\]
die Relationen:
\weq{2}{C \cdot \alpha^2 = 2N \cdot \tau - 2\gamma^2}
\weq{3}{C \cdot \beta^2 = 2P \cdot \tau - 2\delta^2\,.}

223. \textit{Die simultanen Invarianten der vier linearen Covarianten.} Die vier linearen Formen $\alpha$, $\beta$, $\gamma$, $\delta$ liefern ferner sechs simultane Invarianten:
\[
(\beta\alpha),\quad (\gamma\alpha),\quad (\delta\beta);\quad (\delta\alpha),\quad (\delta\gamma),\quad (\gamma\beta)\,.
\]
Die drei ersten derselben haben wir bereits als Hilfsinvarianten $M$, $N$, $P$ eingef\"uhrt; die Formen $(\delta\alpha)$ und $(\gamma\beta)$ stimmen mit $-R$ \"uberein; es bleibt also nur noch die Form $(\delta\gamma)$ zu berechnen. Nun ist:
\weq{1}{(\delta,\ \gamma) = (\tau\beta)(\tau\gamma) = ((\tau\gamma)\tau_x,\ \beta) = ((\tau_1\alpha)(\tau\tau_1)\tau_x,\ \beta)\,.}
Gem\"ass dem Identit\"atssatze ist aber
\[
(\tau_1\alpha)\tau_x - (\tau\alpha)\tau_{1x} = -(\tau\tau_1)\alpha_x\,.
\]
Vertauscht man also in (1) $\tau$ mit $\tau_1$ und nimmt die halbe Summe, so kommt:
\weq{2}{(\delta\gamma) = \left(-\tfrac{1}{2}(\tau\tau_1)^2 \alpha_x,\ \beta_x\right) = \tfrac{1}{2} M \cdot C\,.}
Damit ist auch die sechste dieser simultanen Invarianten auf bekannte Formen reducirt. Es tritt an uns nunmehr die Aufgabe heran, die hiebei benutzten Hilfsinvarianten $M$, $N$, $P$ als ganze Functionen der drei Fundamentalinvarianten darzustellen. Was die letzte dieser drei Formen betrifft, so k\"onnen wir sie unmittelbar auf die beiden andern reduciren. Denn \"uberschiebt man die Relation (1) Nr.\,222 zweimal \"uber $\tau$, so findet man unmittelbar:
\weq{3}{P = \tfrac{1}{2}\left\{2BM - AN\right\}\,.}
Demnach bleibt uns nur mehr die Aufgabe, $M$ und $N$ als Functionen von $A$, $B$ und $C$ darzustellen.


% ===== pdf idx 263, printed 249.0 =====
% !PRINTED 249
\weberpage{249}

224. \textit{Berechnung der Form N.} Das Princip, das uns bei Aufstellung von Relationen leitet, lässt sich dahin formuliren: Ein und dieselbe Form in zweifacher Weise durch Ueberschiebung darzustellen. Nun war $N$ definirt durch

\weq{1}{N = (\tau\alpha)^2 = (jj_1)^2(j_1\alpha)(j\alpha).}

Es handelt sich also nur darum, das symbolische Product $(jj_1)^2(j_1\alpha)(j\alpha)$ durch Ueberschiebung anderer Formen zu gewinnen. Der symbolische Factor $(j\alpha)$ veranlasst uns, die erste Ueberschiebung von $j$ über $\alpha = -(ji)^2j_x$ zu berechnen.

Man erhält:
\[
\begin{aligned}
(j,\alpha) &= -j_x^2(jj_1)(j_1i)^2\\
&= -j_x(jj_1)(j_1i)\{j_{1x}(ji) - i_x(jj_1)\}\\
&= -j_xj_{1x}(jj_1)(j_1i)(ji) + j_xi_x(jj_1)^2(j_1i).
\end{aligned}
\]
Der erste Term rechts ist null; denn er ändert durch Vertauschung von $j$ mit $j_1$ nur sein Zeichen. Der zweite Term ist aber nichts anderes, als $-\vartheta$; denn man hat:
\[
-\vartheta = (\tau,\,i) = \big((jj_1)^2j_xj_{1x},\,i_x^2\big) = j_xi_x(jj_1)^2(j_1i).
\]
Daher hat man zunächst die Relation

\weq{2}{(j,\alpha) = -\vartheta.}

Ueberschiebt man aber diese Gleichung zweimal über sich selbst, so kommt:

\weq{3}{(j\alpha)(j_1\alpha)(jj_1)^2 = (\vartheta,\,\vartheta)^2.}

Es ist aber nach der Theorie der quadratischen Formen

\weq{4}{2(\vartheta,\,\vartheta)^2 = A_{ii}A_{\tau\tau} - A_{i\tau}^2 = AC - B^2.}

und demnach durch Comparation von (1), (3) und (4)
\[
N = \frac{1}{2}(AC-B^2).
\]

225. \textit{Hilfssatz zur Berechnung von M.} Es erübrigt noch, die Form $M$ zu berechnen; zu dem Zwecke beweisen wir zuerst die Hilfsrelation
\[
(f,\,\tau)^2 = -\left\{\frac{2}{3}Aj + i\alpha\right\},
\]
indem wir von einem der symbolischen Producte rechts, etwa von $Aj$, ausgehen und dasselbe geeignet umformen. Es ist
\[
-\frac{2}{3}Aj = -\frac{2}{3}(ab)^4(ai)(bi)\cdot j_x^3.
\]
Gemäss dem Identitätssatze ist
\[
(ab)^3j_x^3 = (aj)^3b^3 - 3(aj)^2(bj)b_x^2a_x + 3(aj)(bj)^2b_xa_x^2 - (bj)^3a_x^3,
\]
%CONT>>


% ===== pdf idx 264, printed 250.0 =====
% !PRINTED 250
\weberpage{250}

%<<CONT
also, da der erste und letzte Term rechts wegen des Reducenten $(aj)^3$ verschwinden:
\[
\begin{aligned}
-\frac{2}{3}Aj &= \left\{2(ab)(aj)^2(bj)b_x^2a_x - 2(ab)(aj)(bj)^2b_xa_x^2\right\}(ai)(bi)\\
&= 2(ab)(aj)(bj)(ai)(bi)b_xa_x\{(aj)b_x - (bj)a_x\}\\
&= 2(ab)^2(aj)(bj)j_x\cdot(ai)(bi)b_xa_x.
\end{aligned}
\]
Unter Benutzung des Productsatzes
\[
(ai)(bi)a_xb_x = \frac{1}{2}\left\{(ai)^2b_x^2 + (bi)^2a_x^2 - (ab)^2i_x^2\right\}
\]
geht das letzte symbolische Product über in
\[
\begin{aligned}
-\frac{2}{3}Aj &= (ab)^2(aj)(bj)(ai)^2b_x^2j_x\\
&\quad + (ab)^2(aj)(bj)(bi)^2a_x^2j_x - (ab)^4(aj)(bj)j_x\cdot i.
\end{aligned}
\]
Die beiden ersten Terme rechts repräsentiren die nämliche Form; der zweite Term zerfällt in $(i,\,j)^2\cdot i = -\alpha\cdot i$. Demnach wird
\[
-\frac{2}{3}Aj = 2(ai)^2(ab)^2(aj)(bj)b_x^2j_x + i\cdot\alpha.
\]
Der symbolische Factor $(ai)^2$ des ersten Termes rechts veranlasst uns, da ja $(ai)^2a_x^3 = -j_x^3$ das Symbol $j_1' = -(ai)^2a_1$, $j_2' = -(ai)^2a_2$ einzuführen. Dadurch wird
\[
\begin{aligned}
-\frac{2}{3}Aj &= -2(j'b)^2(j'j)(bj)b_x^2j_x + i\alpha\\
&= -\{(j'b)j_x - (jb)j_x'\}(j'b)(j'j)(bj)b_x^2 + i\alpha\\
&= -(j'j)^2(j'b)(bj)b_x^3 + i\alpha\\
&= \big(f,\,(jj')^2j_xj_x'\big)^2 + i\alpha,
\end{aligned}
\]
oder, weil ja $\tau = (jj')^2j_xj_x'$:

\weq{1}{(f,\,\tau)^2 = -\left\{\frac{2}{3}Aj + i\alpha\right\}.}

226. \textit{Berechnung der Form M.} Schiebt man nun diese Relation
\[
-3(a\tau)^2a_x^3 = 2Aj + 3i\alpha
\]
dreimal über die Form
\[
(j,\,i) = (ji)j_x^2i_x,
\]
so entsteht links:
\[
-3(a\tau)^2(ji)(aj)^2(ai) = +3(a\tau)(aj)^2(ai)\{(ai)(\tau j) + (ja)(\tau i)\},
\]
d.\,i., weil wegen $(aj)^3a_x^2 = 0$ der zweite Term rechts verschwindet:
\[
\begin{aligned}
&= 3(a\tau)(\tau j)(aj)^2(ai)^2 = -3(ai)^2(aj)^2(a\tau)(j\tau)\\
&= -3\big((ai)^2(aj)^2a_xj_x,\,\tau_x^2\big)^2.
\end{aligned}
\]


% ===== pdf idx 265, printed 251.0 =====
% !PRINTED 251
\weberpage{251}

Nun ist aber $-(ai)^2a_x^3 = j$; daher lässt sich diese Ueberschiebung auch schreiben

\weq{1}{= +3\big((jj_1)^2j_xj_{1x},\,\tau\big)^2 = 3(\tau,\,\tau)^2 = 3C.}

Auf der rechten Seite erhalten wir durch die dreimalige Ueberschiebung: 1) das Glied

\weq{2}{(2Aj,\,(j,\,i))^3 = 2A\cdot(j_1i)(jj_1)^2(ji) = 2A\cdot(\tau,\,i)^2 = 2A\cdot B,}

2) das Glied
\[
(3i_1\cdot\alpha,\,(ji)j_x^2i_x)^3 = (ji)(i_1j)^2(\alpha i) + 2(ji)(\alpha j)(i_1j)(i_1i).
\]
Der zweite Term rechts ist null; denn er ändert durch Vertauschung von $i$ mit $i_1$ nur sein Zeichen. Der erste Term kann aber dargestellt werden durch

\weq{3}{\begin{aligned}
(ji_1)^2(ji)(\alpha i) &= \big((ji_1)^2j_x,\,(\alpha i)i_x\big)\\
&= (-\alpha_x,\,(\alpha i)i_x)\\
&= -(\alpha i)^2 = -M.
\end{aligned}}

Die gesuchte Relation wird demnach
\[
3C = 2AB - M.
\]
Fassen wir die bisher gewonnenen Resultate zusammen, so erhalten wir:

\weq{4}{(\tau\alpha)^2 = N = \frac{1}{2}(AC-B^2)}

\weq{5}{(i\alpha)^2 = M = 2AB - 3C}

\weq{6}{(\tau\beta)^2 = P = \frac{1}{4}(9AB^2 - A^2C - 12BC).}

Die Gleichnng (5) lehrt auch, dass wir berechtigt waren, $(\tau,\,\tau)^2 = C$ an Stelle von $(f^2,\,i^5)^{10}$ einzuführen. % [sic] \textquotedbl{}Gleichnng\textquotedbl{} printed for \textquotedbl{}Gleichung\textquotedbl{}
Denn $(i\alpha)^2$ ist das Glied
\[
G = \big((ai_1)^2(ai_2)^2a_x\cdot(bi_3)^2(bi_4)^2b_x,\,i_{5x}^2\big)^2
\]
dieser Ueberschiebung $(f^2,\,i^5)^{10}$.

227. \textit{Die Ueberschiebungen der quadratischen Formen $i,\,\tau,\,\vartheta$.} Eine weitere Reihe von Relationen, deren wir später bedürfen, liefern die gegenseitigen Ueberschiebungen der quadratischen Formen des Systemes, sowie die Ueberschiebungen dieser Formen über die vier linearen Formen. Sie sind leicht zu ermitteln, da die simultanen Systeme dieser Formen niedrigern Grades nach den früheren Entwicklungen bekannt sind. Es ist überhaupt, wie wir schon bisher sehen konnten, nützlich, das Studium der Systeme höherer Formen dadurch zu vereinfachen, dass man insbesondere die Ueberschiebungen der linearen und quadratischen Formen des Systemes in den Vordergrund stellt.

Die drei quadratischen Formen des Systemes waren $i,\,\tau,\,\vartheta$.


% ===== pdf idx 266, printed 252.0 =====
% !PRINTED 252
\weberpage{252}

Da $\vartheta$ Functionaldeterminante der beiden ersten Formen ist, so haben wir nach der Theorie der quadratischen Formen unmittelbar folgende Relationen (vgl. Nr. 127):

\weq{1}{(i,\,\vartheta) = \frac{1}{2}\left\{\cdot B\cdot i - A\cdot\tau\right\}} % [sic] stray dot before B as printed

\weq{2}{(\tau,\,\vartheta) = \frac{1}{2}\left\{C\cdot i - B\cdot\tau\right\}}

\weq{2a}{\left.\begin{aligned}
(i,\,\vartheta)^2 &= A_{i\vartheta} = 0\\
(\tau,\,\vartheta)^2 &= A_{\tau\vartheta} = 0
\end{aligned}\right\}}

\weq{3}{-2\vartheta^2 = A\tau^2 - 2Bi\tau + Ci^2.}

Aus der letzten dieser Relationen geht ohne weiteres die Beziehung hervor, welche zwischen dem Quadrat der schiefen Invariante $R$ und den übrigen fundamentalen Invarianten besteht. Denn es war
\[
R = (\vartheta,\,\alpha^2)^2.
\]
Ersetzt man also in dieser Relation $x$ durch $\alpha$, so kommt:

\weq{4}{-2R^2 = AN^2 - 2BMN + CM^2.}

228. \textit{Die Ueberschiebungen von $\vartheta$ über $\alpha,\,\beta,\,\gamma,\,\delta$.} Von den Ueberschiebungen der quadratischen Formen $i,\,\tau,\,\vartheta$ will ich, soweit wir sie nicht ohnehin schon berechnet haben, noch jene Formen ermitteln, welche sich bei einmaliger Ueberschiebung der Form $\vartheta$ über $\alpha,\,\beta,\,\gamma,\,\delta$ ergeben. Alle weiteren Ueberschiebungen mögen der selbstständigen Uebung zugewiesen sein.

Um nun aber $(\vartheta,\,\alpha),\,(\vartheta,\,\beta),\,(\vartheta,\,\gamma),\,(\vartheta,\,\delta)$ zu berechnen, benutzen wir einfach die schon bei Gelegenheit der quadratischen Formen gegebene Reihenentwicklung [vgl. Nr. 125 (2)]:
\[
\vartheta_x\vartheta_y = (i\tau)i_x\tau_y - \frac{B}{2}(xy) = (i\tau)i_y\tau_x - \frac{B}{2}(yx).
\]
Wir ersetzen in diesen Gleichungen $x$ resp. $y$ der Reihe nach durch $\alpha,\,\beta,\,\gamma,\,\delta$ und erhalten nach einfacher Umformung:

\weq{1}{(\vartheta,\,\alpha) = (i\tau)(i\alpha)\tau_x - \frac{B}{2}\alpha = -(\tau\beta)\tau_x - \frac{B}{2}\alpha = -\delta - \frac{B}{2}\alpha}

\weq{2}{\begin{aligned}
(\vartheta,\,\beta) &= (i\tau)(i\beta)\tau_x - \frac{B}{2}\beta = (i\tau)(ii_1)(i_1\alpha)\tau_x - \frac{B}{2}\beta\\
&= \frac{A}{2}(\tau\alpha)\tau_x - \frac{B}{2}\beta = +\frac{A}{2}\gamma - \frac{B}{2}\beta
\end{aligned}}

\weq{3}{(\vartheta,\,\gamma) = (i\tau)(\tau\gamma)i_x + \frac{B}{2}\gamma = (i\tau)(\tau\tau_1)(\tau_1\alpha)i_x + \frac{B}{2}\gamma = -\frac{C}{2}\beta + \frac{B}{2}\gamma}

\weq{4}{(\vartheta,\,\delta) = (i\tau)(\tau\delta)i_x + \frac{B}{2}\delta = (i\tau)(\tau\tau_1)(\tau_1\beta)i_x + \frac{B}{2}\delta = +\frac{AC}{4}\alpha + \frac{B}{2}\delta.}


% ===== pdf idx 267, printed 253.0 =====
% !PRINTED 253
\weberpage{253}

In ähnlicher Weise erhält man unter Benutzung dieser vier Relationen
die Ueberschiebungen von $\vartheta$ über Producte der linearen Covarianten.

Damit will ich diese Untersuchungen vorläufig abschliessen. Wir
werden im folgenden Paragraphen noch eine Reihe anderer Ueberschiebungen
berechnen, so weit wir derselben zur Lösung der dort
gestellten Aufgabe bedürfen. Insbesondere werden wir an geeigneter
Stelle die Ueberschiebungen von $f$ über die linearen Covarianten $\alpha$, $\beta$,
$\gamma$, $\delta$ durch Formen des Systemes und durch Ueberschiebungen derselben
darstellen.

\begin{center}
\S~24.

\textbf{Typische Darstellung der Form fünften Grades.}
\end{center}

229. \textit{Begriff der typischen Darstellung.} Unter typischer Darstellung
einer Form $f = a_x^n$ versteht man im Allgemeinen eine derartige
Darstellung derselben, dass die Coefficienten Invarianten und
die beiden Variabeln irgend zwei Covarianten des Systemes sind. Der
neue Ausdruck für $f$ ist alsdann der unveränderliche Typus für alle
Formen, die aus $f$ durch lineare Transformation hervorgehen; und umgekehrt:
Haben zwei Formen $f$ die nämliche typische Form, so können
sie stets durch lineare Transformation in einander übergeführt werden.

Für Formen ungeraden Grades kann die typische Darstellung unter
anderm auf folgendem Wege erreicht werden. Dieselben besitzen im
allgemeinen Falle und wenn $n > 3$ stets mindestens zwei lineare
Covarianten $\alpha_x$ und $\beta_x$ mit nicht verschwindender Determinante $(\alpha\beta)$,
wie \gesperrt{Clebsch} (vgl. \glqq Binäre Formen\grqq{} \S~90, Seite 357) gezeigt hat, und
diese werden als neue Variable für die typische Darstellung benutzt.
Erhebt man nämlich die Identität:
\[
a_x(\alpha\beta) = \alpha_x(\alpha\beta) - \beta_x(a\alpha)
\]
auf die $n^{\mathrm{te}}$ Potenz, so kommt:
\[
a_x^n \cdot (\alpha\beta)^n = \alpha_x^n(\alpha\beta)^n - \binom{n}{1}\alpha_x^{n-1}\beta_x(\alpha\beta)^{n-1}(a\alpha) + \cdots \pm \beta_x^n(a\alpha)^n
\]
oder:
\weq{I}{(\alpha\beta)^n \cdot f = (f,\ \beta^n)^n \cdot \alpha_x^n - \binom{n}{1}(f,\ \beta^{n-1}\alpha)^n \cdot \alpha_x^{n-1}\beta_x
+ \binom{n}{2}(f,\ \beta^{n-2}\alpha^2)^n \cdot \alpha_x^{n-2}\beta_x^2 - \cdots}

Hierin sind aber die Variabeln Covarianten, die Coefficienten Invarianten,
und man hat nur noch die Aufgabe zu lösen, diese Invarianten:
\weq{II}{(f,\ \beta^n)^n,\quad (f,\ \beta^{n-1}\alpha)^n,\quad (f,\ \beta^{n-2}\alpha^2)^n \cdots (f,\ \alpha^n)^n,\quad (\alpha\beta)^n}
durch die fundamentalen Invarianten des Systemes zu ersetzen\ldots{} Man
%CONT>>


% ===== pdf idx 268, printed 254.0 =====
% !PRINTED 254
\weberpage{254}

%<<CONT
hätte auch, wie dies \gesperrt{Hermite} gethan hat, statt der linearen Formen
$\alpha$ und $\beta$ die linearen Factoren einer quadratischen Covariante des
Systemes benutzen können. Dadurch ist man freilich genöthigt, Irrationalitäten
einzuführen; aber die typischen Coefficienten sind alsdann
vom niedrigeren Grade in den Coefficienten der ursprünglichen Form.

Um nun die Form $f = a_x^5$ typisch darzustellen, benutzen wir zunächst
die beiden linearen Covarianten
\[
\alpha_x = (f,\ i^2)^4, \quad \beta_x = -(f,\ i^3)^5.
\]
Wir erhalten alsdann, weil $(\alpha\beta) = -M = 3C - 2AB$:
\[
(-M)^5 \cdot f = \alpha_x^5 \cdot (f,\ \beta^5)^5 - 5\alpha_x^4\beta_x \cdot (f,\ \beta^4\alpha)^5
+ 10\alpha_x^3\beta_x^2(f,\ \beta^3\alpha^2)^5 - \cdots - \beta_x^5(f,\ \alpha^5)^5
\]
und unsere Aufgabe besteht nun darin, die sechs Invarianten
\[
(f,\ \beta^5)^5,\quad (f,\ \beta^4\alpha)^5,\quad (f,\ \beta^3\alpha^2)^5,\quad (f,\ \beta^2\alpha^3)^5,\quad (f,\ \beta\alpha^4)^5,\quad (f,\ \alpha^5)^5
\]
durch die Invarianten $A$, $B$, $C$, $R$ des Systemes darzustellen.

230. \textit{Berechnung der Ueberschiebungen $(f, \alpha^{\varrho})^{\varrho}$.} Zu dem Zwecke
haben wir zunächst die fünf Ueberschiebungen von $f$ über $\alpha$ durch
andere Ueberschiebungen darzustellen, eine Aufgabe, die wir hier,
wenn auch nicht in voller Ausführlichkeit, erledigen wollen. Man
findet zunächst:
\weq{I}{(f,\ \alpha) = 2\left((f,\ j)^2,\ i\right)}
\weq{II}{(f,\ \alpha^2)^2 = \frac{3}{2}\,\alpha\tau - Ai\alpha + \lambda j,}
wobei $\lambda$ den Werth $B - \frac{2}{3}A^2$ besitzt. Um zur ersten Gleichung zu
gelangen, hat man in $(a\alpha)a_x^4$ die Form $\alpha_x$ durch ihren Werth
$-(ji)^2 j_x$ zu ersetzen und das so entstehende symbolische Product durch
den Identitätssatz umzuformen. Indem man dann anderntheils die Polare
%CONT>>


% ===== pdf idx 269, printed 255.0 =====
% !PRINTED 255
\weberpage{255}

%<<CONT
$(f,\ j)_y^2$ berechnet und darin $y$ durch $i$ ersetzt, kommt man durch abermalige
Anwendung derselben Identität direct zu Gleichung (I).

Complicirter gestaltet sich die Herstellung der zweiten Relation.
Man benutzt die eben gewonnene Relation (I) und überschiebt sie
nochmals über die Form $\alpha$, indem man der Einfachheit halber für
$(f,\ j)^2$ zunächst die Bezeichnung $p_x^4$ einführt. Dies liefert
\[
4(f,\ \alpha^2)^2 = 6p_x^2 i_x (p\,i)(p\alpha) i_x + 2p_x^3(p\,i)(i\alpha),
\]
oder, indem man im zweiten Terme rechts $p_x(i\alpha)$ durch $i_x(p\alpha) - \alpha_x(p\,i)$
ersetzt:
\weq{1}{4(f,\ \alpha^2)^2 = 8(p\,i)(p\alpha)p_x^2 i_x - 2(p\,i)^2 \alpha_x p_x^2.}

Man hat nun zu berücksichtigen, dass die beiden Glieder der schon
mehrmals benutzten Polare $(f,\ j)_y^2$ einander gleich sind, da ihre Differenz
wegen des sich einstellenden Factors $(aj)^3$ verschwindet, dass also
\weq{2}{p_x^3 p_y = (aj)^2 a_x^3 j_y}
\weq{3}{p_x^3 p_y = (aj) a_x^2 j_x a_y.}
Mit Hilfe dieser beiden Gleichungen lässt sich nun leicht in (1) das
Symbol $p$ wiederum eliminiren, da auch die beiden Glieder von $p_x^2 p_y^2$
in (3) identisch gleich sind, und man findet so:
\weq{4}{4(f,\ \alpha^2)^2 = 8(aj)^2(ai)(j\alpha)a_x^2 i_x - 2(aj)^2(ai)^2 a_x j_x \alpha_x \cdot}

Hierin ist der zweite Term incl. Vorzeichen nichts anderes als $+2\tau \cdot \alpha$.
Der erste Term aber geht durch die Substitution $i_x(aj) = j_x(ai) - a_x(j\,i)$
über in
\weq{5}{8(aj)^2(ai)(j\alpha)a_x^2 i_x = 8(aj)(ai)^2(j\alpha)j_x a_x^2 - 8(aj)(ai)(j\alpha)(j\,i)a_x^3,}
wobei [vgl. auch Nr. 221, und Nr. 224 (2)]
\[
(aj)(ai)^2(j\alpha)j_x a_x^2 = \left(-(ai)^2 a_x^3,\ -(j\alpha)j_x^2\right) = (j,\ \vartheta) = \frac{\tau \cdot \alpha}{2}
\]
\[
-(aj)(ai)(j\alpha)(j\,i)a_x^3 = \left(a_x^5,\ -(j\alpha)(j\,i)i_x j_x\right)^2 = (f,\ (\vartheta,\ i))^2.
\]
Demnach wird
\[
4(f,\ \alpha^2)^2 = 6\alpha\tau + 8(f,\ (\vartheta,\ i))^2.
\]
Führen wir hierin nun den in Nr. 227 (1) berechneten Werth von
$(\vartheta,\ i)$ ein und eliminiren die sich hiebei einstellende Form $(f,\ \tau)^2$ mittelst
Gleichung (1) Nr. 225, so ergiebt sich in der That die oben gegebene
Relation (II). Aus den beiden Relationen (I) und (II) gehen
nun die übrigen Ueberschiebungen direct hervor. Man überschiebt (II)
abermals über $\alpha$ und findet:
\weq{III}{(f,\ \alpha^3)^3 = \alpha \cdot \gamma - \frac{2}{3}A \cdot \alpha \cdot \beta - \lambda \cdot \vartheta}
%CONT>>


% ===== pdf idx 270, printed 256.0 =====
% !PRINTED 256
\weberpage{256}

%<<CONT
und hieraus durch weitere Ueberschiebungen
\weq{IV}{(f,\ \alpha^4)^4 = \mu \cdot \alpha + \lambda \cdot \delta}
\weq{V}{(f,\ \alpha^5)^5 = -\lambda \cdot R,}
wobei $\mu$ den Werth $\frac{1}{2}N - \frac{1}{3}AM + \lambda\frac{B}{2}$ besitzt.

231. \textit{Berechnung der typischen Coefficienten.} Die soeben berechneten
Ueberschiebungen gestatten uns nun in Verbindung mit den in
Nr. 222 gegebenen Relationen für die Quadrate der Functionaldeterminanten
$\beta^2$, $\gamma^2$, $\delta^2$ in einfacher Weise nicht nur die typischen Coefficienten,
sondern überhaupt alle Ueberschiebungen von $f$ über ein Product
$\alpha^{\varkappa}\beta^{\lambda}\gamma^{\mu}\delta^{\nu}$, $\varkappa + \lambda + \mu + \nu = 5$, durch die Invarianten des
Systemes darzustellen. Insbesondere ist:
\[
(f,\ \alpha^5)^5 = -\lambda \cdot R
\]
\[
(f,\ \alpha^4\beta)^5 = \left((f,\ \alpha^4)^4,\ \beta\right) = -\mu \cdot M + \lambda P
\]
\[
(f,\ \alpha^3\beta^2)^5 = \left((f,\ \alpha^3)^3,\ \beta^2\right)^2 = (f,\ \alpha^3 \cdot i)^5 M - (f,\ \alpha^5)^5 \frac{A}{2}
\]
\[
{} = +R\left(M + \frac{\lambda}{2}A\right), \quad \text{[weil: } (f,i)^2 = -j,\ -(j,\alpha)^3 = (\vartheta,\alpha^2)^2 = R\text{]}
\]
\[
(f,\ \alpha^2\beta^3)^5 = \left((f,\ \alpha^2)^2,\ \beta\left(Mi - \frac{A\alpha^2}{2}\right)\right)^3
\]
\[
{} = (f,\ \alpha^2\beta i)^5 M - (f,\ \alpha^4\beta)^5 \frac{A}{2}, \quad \text{[weil: } (-j,\alpha^2\beta)^3 = (\vartheta,\alpha\beta)^2\text{]}
\]
\[
{} = M\left(\frac{A}{2}N - \frac{B}{2}M\right) - \frac{A}{2}(\lambda P - \mu M), \quad \text{[vgl. Nr. 228 (1)]}
\]
\[
(f,\ \alpha \cdot \beta^4)^5 = -AMR - \frac{A^2}{4} \cdot \lambda \cdot R
\]
\[
(f,\ \beta^5)^5 = (f,\ \beta i^2)^5 M^2 - (f,\ \beta\alpha^2 i)^5 A \cdot M + (f,\ \beta\alpha^4)^5 \frac{A^2}{4}, \quad [-(j,i)^2 = \alpha]
\]
\[
{} = -M^3 + AM\left(\frac{A}{2}N - \frac{B}{2}M\right) + \frac{A^2}{4}(\lambda P - \mu M).
\]

232. \textit{Typische Darstellung der Covariante $j$.} Ehe wir die typischen
Darstellungen der Form $f$ weiter verfolgen, wollen wir auch für die
Covariante $j$ eine solche geben, indem wir hiezu die linearen Covarianten
$\alpha$ und $\delta$ benutzen. Nach dem binomischen Lehrsatze erhält
man:
\[
(\alpha\delta)^3 j_x^3 = (j\delta)^3 \alpha_x^3 - 3(j\delta)^2(j\alpha)\alpha_x^2\delta_x + 3(j\delta)(j\alpha)^2\alpha_x\delta_x - (j\alpha)^3\delta_x^3.
\]
Hierin ist:

1) $(j\alpha)^3 = ((j,\alpha),\alpha^2)^2 = (-\vartheta,\alpha^2)^2 = -R$ (vgl. Nr. 224 (2), 220 (8)),

2) $(j,\alpha^2\delta)^3 = ((j,\alpha),\alpha\delta)^2 = (-\vartheta,\alpha\delta)^2 = (-(\vartheta,\alpha),\delta)$
\[
{} = \left(\delta + \frac{B}{2}\alpha,\ \delta\right) = +\frac{B}{2} \cdot R \quad \text{(vgl. Nr. 228 (1)),}
\]

3) $(j,\alpha\delta^2)^3 = (-\vartheta,\delta^2)^2 = \left(-\vartheta,\ P\tau - C \cdot \frac{\beta^2}{2}\right)^2 \quad \text{(vgl. Nr. 222 (3)),}$
\[
{} = -P \cdot (\vartheta,\tau)^2 + \frac{C}{2}(\vartheta,\beta^2)^2,
\]
oder, weil nach Nr. 227 (2a) die Form $(\vartheta,\tau)^2 = 0$, so kommt:
\[
(j,\ \alpha\delta^2)^3 = \frac{C}{2}\left((\vartheta,\beta),\beta\right) = \frac{C}{2}\left(\frac{A}{2}\gamma - \frac{B}{2}\beta,\ \beta\right) = -\frac{CA}{4}R,
\]

4) $(j,\delta^3)^3 = \left(j,\ \delta\left(P\tau - \frac{C}{2}\beta^2\right)\right)^3 = P \cdot (j,\delta\tau)^3 - \frac{C}{2}(j,\delta\beta^2)^3,$

oder, weil $(j,\tau)^2$ gemäss der Theorie cubischer Formen verschwindet:
\[
(j,\ \delta^3)^3 = -\frac{C}{2}\left(j,\ \delta\left(Mi - \frac{A\alpha^2}{2}\right)\right)^3 = +\frac{CR}{2}\left(M + \frac{AB}{4}\right).
\]


% ===== pdf idx 271, printed 257.0 =====
% !PRINTED 257
\weberpage{257}
Die typische Darstellung von $j$ wird also, da $(\alpha\delta) = +\mathbf{R}$, nach
Division mit $\mathbf{R}$

\weq{I}{
\mathbf{R}^2\cdot j = \delta^3 + \frac{3\mathbf{B}}{2}\cdot\delta^2\alpha + \frac34\,\mathbf{A}\mathbf{C}\cdot\delta\alpha^2 + \frac{\mathbf{C}}{2}\left(\mathbf{M}+\frac{\mathbf{A}\mathbf{B}}{4}\right)\cdot\alpha^3.
}

Nehmen wir an, die cubische Form $j$ besitze, gleich null gesetzt, drei
verschiedene Wurzeln $m_1\,m_2\,m_3$, was ja im Allgemeinen auch der
Fall sein wird, da sonst die Invariante $\mathbf{C} = (\tau,\tau)^2$ des Systems
verschwinden würde, so können wir schreiben:

\weq{II}{
\mathbf{R}^2\cdot j = (\delta - m_1\alpha)(\delta - m_2\alpha)(\delta - m_3\alpha).
}

233. \textit{Darstellung der Form $f$ durch drei fünfte Potenzen.} Eine
der wichtigsten typischen Darstellungen der Form $f$ ist die typische
Darstellung durch fünfte Potenzen. Der Umstand, dass

\weq{1}{
(f,\,j)^3 = 0,
}

zeigt, dass es immer möglich ist, $f$ in die Gestalt

\[
f = k_1\,p_x^5 + k_2\,q_x^5 + k_3\,r_x^5
\]

zu bringen, wo $p_x$, $q_x$, $r_x$ die drei linearen als von einander
verschieden angenommenen Factoren der cubischen Form $j$ sind. Denn
nach den Ueberschiebungsgesetzen (vergl. Nr. 38) ist stets:

\[
(k_1p_x^5 + k_2q_x^5 + k_3r_x^5,\ p_xq_xr_x)^3 = 0,
\]

und diese Gleichung ist ja gerade die Definitionsgleichung der Co\-
variante $j = p_xq_xr_x$ (vgl. Nr. 215).

Legen wir $j$ in der typischen Darstellung (II) zu Grunde, dann ist

\[
p = \delta - m_1\alpha
\]
\[
q = \delta - m_2\alpha
\]
\[
r = \delta - m_3\alpha,
\]

und demnach die typische Form von $f$

\weq{2}{
f = k_1(\delta - m_1\alpha)^5 + k_2(\delta - m_2\alpha)^5 + k_3(\delta - m_3\alpha)^5.
}

Die Berechnung der Constanten $k_1\,k_2\,k_3$ kann in folgender Weise
vorgenommen werden.

Wir überschieben Gleichung (2) fünfmal über $\alpha^5$, resp. $\alpha^4\delta$,
$\alpha^3\delta^2$ und erhalten:

\weq{3}{
\left.
\begin{aligned}
(f,\ \alpha^5)^5 &= -\ (k_1 + k_2 + k_3)\ \mathbf{R}^5\\
(f,\ \alpha^4\delta)^5 &= -\ 5(k_1m_1 + k_2m_2 + k_3m_3)\ \mathbf{R}^5\\
(f,\ \alpha^3\delta^2)^5 &= -\ 10(k_1m_1^2 + k_2m_2^2 + k_3m_3^3)\ \mathbf{R}^5
\end{aligned}
\right\}
}
% [sic] letzter Exponent m_3^3 so gedruckt (statt m_3^2, vgl. die beiden vorigen Glieder)

Die Werthe der linken Seite dieser drei Gleichungen ergeben sich
unmittelbar aus den Gleichungen (III), (IV), (V) Nr. 230, und da die
Determinante dieses Systemes (3) von linearen Gleichungen nicht
%CONT>>


% ===== pdf idx 272, printed 258.0 =====
% !PRINTED 258
\weberpage{258}
%<<CONT
verschwindet (wegen $(\tau_1\tau)^2 \gtrless 0$), so lassen sich aus ihm die drei
Constanten $k_1$, $k_2$, $k_3$ eindeutig berechnen.

234. \textit{Typische Darstellung von $f$, wenn $\mathbf{C} = 0$.} Sind nun aber die
drei Wurzeln der Covariante $j = 0$ nicht von einander verschieden,
sind vielmehr zwei derselben einander gleich, dann ist die Darstellung
durch fünfte Potenzen allein nicht mehr möglich. Immer noch aber
liefert die typische Darstellung von $f$ unter Zugrundelegung der Linear\-
factoren von $j$ ein sehr einfaches Resultat, wovon man sich sofort
überzeugen kann. In diesem Falle ist nämlich:

\weq{1}{
j = p_x^2\,q_x
}

und demnach

\weq{2}{
(f,\,j)^3 = (ap)^2(aq)\,a_x^2 = 0.
}

Erhebt man also die Identität

\[
a_x(pq) = p_x(aq) - q_x(ap)
\]

auf die fünfte Potenz, so verschwinden in der binomischen Entwicklung
rechts alle Glieder, in denen die Klammerfactoren $(ap)^2(aq)$ auftreten.
Demnach reducirt sich dieselbe auf

\weq{3}{
(pq)^5\cdot f = (aq)^5p_x^5 - 5(ap)(aq)^4q_xp_x^4 - (ap)^5q_x^5,
}

oder also auf

\weq{4}{
(pq)^5\cdot f = c_1p_x^5 + 5c_2p_x^4q_x + c_3q_x^5,
}

wo $c_1\,c_2\,c_3$ gewisse Constante sind, die wir nun berechnen wollen.

235. \textit{Berechnung der Coefficienten der Gleichung (4).} Zu dem
Zwecke gehen wir wiederum von der typischen Darstellung der Co\-
variante $j$ aus. [Vergl. Nr. 232, (I).] Da aber $j$ nach Voraussetzung
einen linearen Doppelfactor besitzt, und demnach ihre Discriminante
$\mathbf{C} = (\tau,\tau)^2 = 0$ ist, so reducirt sich dieselbe auf

\[
\mathbf{R}^2\cdot j = \delta^2\left(\delta + \frac32\,\mathbf{B}\alpha\right).
\]

Setzen wir also

\weq{5}{
\varrho = \delta + \frac32\,\mathbf{B}\alpha,
}

so hat man:

\weq{5a}{
\mathbf{R}^2\cdot j = \delta^2\cdot\varrho,
}

und

\weq{5b}{
(\delta\varrho) = \frac32\,\mathbf{B}(\delta\alpha) = -\,\frac32\,\mathbf{B}\cdot\mathbf{R}.
}

Die Gleichung (4) geht also über in

\[
(\delta\varrho)^5\cdot f = (a\varrho)^5\delta_x^5 - 5(a\delta)(a\varrho)^4\delta_x^4\varrho_x - (a\delta)^5\varrho_x^5,
\]

oder

\weq{6}{
(\delta\varrho)^5f = c_1\delta^5 + 5c_2\delta^4\varrho + c_3\varrho^5.
}


% ===== pdf idx 273, printed 259.0 =====
% !PRINTED 259
\weberpage{259}
Die Berechnung der Coefficienten $c_i$ bewerkstelligt man nun genau
in derselben Weise wie im allgemeineren Falle Nr. 233. Wir bilden
die fünften Ueberschiebungen dieser Gleichung (6) über $\alpha^5$, $\alpha^4\delta$, $\alpha^3\delta^2$,
und erhalten dadurch drei Gleichungen zur Bestimmung der Grössen $c_i$.

Es ergiebt sich nämlich dadurch, wenn wir einen Augenblick den
Factor $(\delta\varrho)^5$ mit $m\cdot\mathbf{R}^5$ bezeichnen [siehe Gleichung (5b)]:

\weq{6a}{
\left.
\begin{aligned}
\mathbf{R}^5\cdot m\cdot(a\alpha)^5 &= c_1(\delta\alpha)^5 + 5c_2(\delta\alpha)^4(\varrho\alpha) + c_3(\varrho\alpha)^5\\
\mathbf{R}^5\cdot m\cdot(a\alpha)^4(a\delta) &= 5c_2(\delta\alpha)^4(\varrho\delta) + c_3(\varrho\alpha)^4(\varrho\delta)\\
\mathbf{R}^5\cdot m\cdot(a\alpha)^3(a\delta)^2 &= c_3(\varrho\alpha)^3(\varrho\delta)^2
\end{aligned}
\right\}
}

Nun ist gemäss den Relationen (IV) und (V) in Nr. 230

\weq{7}{
(a\alpha)^5 = -\,\lambda\mathbf{R},
}

\[
(a\alpha)^4(a\delta) = \ \mu\mathbf{R}.
\]

Um aber $(f,\,\alpha^3\delta^2)^5$ zu berechnen, benutzen wir am einfachsten
die Relation $\delta^2 = \mathbf{P}\tau - \dfrac{\mathbf{C}}{2}\beta^2$ (vergl. Nr. 222), oder weil $\mathbf{C} = 0$, die
Relation:

\[
\delta^2 = \mathbf{P}\tau.
\]

Dann wird:

\[
\begin{aligned}
(a\alpha)^3(a\delta)^2 &= (f,\,\alpha^3\mathbf{P}\tau)^5 = \mathbf{P}((f,\tau)^2,\,\alpha^3)^3,\\
&= \mathbf{P}\left(\left(-i\alpha - 2\,\frac{\mathbf{A}}{3}\,j\right),\ \alpha^3\right)^3 \qquad \text{[vergl. Nr. 225, (1)]}\\
&= \frac23\,\mathbf{P}\mathbf{A}(-j,\,\alpha^3)^3 = \frac23\,\mathbf{P}\mathbf{A}(-(j,\alpha),\,\alpha^2)^2 \qquad \text{[Nr. 224, (2)]}\\
&= \frac23\,\mathbf{A}\mathbf{P}(\vartheta,\,\alpha)^2 = \frac23\,\mathbf{A}\mathbf{P}\mathbf{R},
\end{aligned}
\]

oder, weil für $\mathbf{C} = 0$ die Invariante $\mathbf{P} = \dfrac94\,\mathbf{A}\mathbf{B}^2$ wird [vgl. Nr. 226, (6)]:

\weq{9}{
(a\alpha)^3(a\delta)^2 = \frac32\,\mathbf{A}^2\mathbf{B}^2\mathbf{R}.
}

Anderntheils ist:

\[
(\delta\alpha) = -\,\mathbf{R}, \quad (\varrho\alpha) = -\,\mathbf{R}, \quad (\varrho\delta) = +\,\frac32\,\mathbf{B}\cdot\mathbf{R}.
\]

Demnach ergeben sich durch Eintragen dieser Werthe aus den
Gleichungen (6a) die folgenden Gleichungen:

\weq{10}{
\left.
\begin{aligned}
m\cdot\lambda\cdot\mathbf{R} &= c_1 + 5c_2 + c_3\\
m\cdot\mu\cdot\mathbf{R} &= {}+\,\frac{15}{2}\,\mathbf{B}\cdot c_2 + \frac32\,\mathbf{B}\cdot c_3\\
-\,\frac32\,m\cdot\mathbf{A}^2\cdot\mathbf{B}^2\cdot\mathbf{R} &= {}+\,\frac94\,\mathbf{B}^2\cdot c_3
\end{aligned}
\right\}
}

Hieraus findet sich zunächst

\[
c_3 = -\,\frac23\,\mathbf{A}^2\cdot m\cdot\mathbf{R}.
\]


% ===== pdf idx 274, printed 260.0 =====
% !PRINTED 260
\weberpage{260}
Indem man alsdann $\mu$ und $\lambda$ durch ihre Werthe (Nr. 230) in $\mathbf{A}$
und $\mathbf{B}$ ersetzt, findet man ferner

\[
c_2 = c_1 = \frac{\mathbf{B}}{6}\cdot m\cdot\mathbf{R}.
\]

Die typische Darstellung der Form $f$ wird demnach für $\mathbf{C} = 0$

\weq{11}{
6\mathbf{R}^4\cdot f = \mathbf{B}\delta^5 + 5\mathbf{B}\delta^4\varrho - 4\mathbf{A}^2\varrho^5.
}

Setzen wir $f = 0$, so haben wir in dieser Darstellung die
\glqq Bring'sche\grqq{} Normalform der Gleichung fünften Grades. Durch die
Substitutionen

\[
\frac{5\mathbf{B}}{-4\mathbf{A}^2} = \mathbf{Z}, \quad \frac{\varrho}{\delta} = t\sqrt[4]{\mathbf{Z}}
\]

geht sie in die von Hermite benutzte Gestalt über

\[
t^5 - t - \frac{1}{5\sqrt[4]{\mathbf{Z}}} = 0.
\]

Sobald also $\mathbf{C} = 0$, lässt sich immer eine lineare Transformation
angeben, wodurch die allgemeine Gleichung fünften Grades in diese
Normalform mit dem einzigen Parameter $X = \dfrac{1}{5\sqrt[4]{\mathbf{Z}}}$ übergeht. Da\-
gegen kann die allgemeine Gleichung fünften Grades nur durch eine
Transformation vierten Grades in diese Normalform gebracht werden.

236. \textit{Bedeutung der Relation $\mathbf{C} = 0$.} Es ist interessant, die Be\-
deutung der Relation $\mathbf{C} = 0$ für das Formensystem noch weiter zu
untersuchen. Man findet alsdann insbesondere zweierlei: 1) Die Co\-
varianten $j$ und $i$ besitzen, sobald $\mathbf{C} = 0$, einen gemeinsamen linearen
Factor, oder mit andern Worten: die Invariante $\mathbf{C}$ ist bis auf einen
Zahlenfactor gleich der Resultante von $j$ und $i$; 2) dieser gemeinsame
Factor von $i$ und $j$ ist nichts anderes als die oben eingeführte lineare
Covariante $\varrho$.

Indem wir die diesbezüglichen Untersuchungen hier anfügen, ver\-
folgen wir einen doppelten Zweck: einmal an einem Beispiel zu zeigen,
wie durch symbolische Rechnung die Resultante zweier Covarianten
abgeleitet wird, dann aber auch, um durch die sich ergebenden Re\-
sultate ein neues Licht auf das Wesen der Bring'schen Normalform
zu werfen.

Wir beginnen damit, die Resultante von $j$ und $i$ zu berechnen.

Angenommen, die beiden linearen Factoren der quadratischen
Form $i$ seien

\[
i = r_x\,s_x;
\]

sie besitze also die Wurzeln $\xi_1 = \dfrac{-r_2}{r_1} = r$, $\xi_2 = \dfrac{-s_2}{s_1} = s$; dann ist,
%CONT>>


% ===== pdf idx 275, printed 261.0 =====
